% GENERATED FILE. Edit canonical lessons, then run npm run generate:books.
% canonical-source-hash: fnv1a64:1538bfdb28dd8d00
% canonical-lessons: TA-C411-kolaru, TA-C411-ceyalila, TA-C411-nacunku, TA-C411-uracu, TA-C411-muttu

\chapter{When the Car Breaks Down}
\label{ch:ta-when-the-car-breaks-down}
\hlchaptermodality{411}

\begin{chapteropening}
\textbf{By the end of this chapter:} \emph{I can talk about a fault and an engine that stops working, and dents, scrapes and bumps.}

Complete the last lesson of chapter 411: I can talk about a fault and an engine that stops working, and dents, scrapes and bumps.
\end{chapteropening}

\section[\texorpdfstring{\ta{கோளாறு} (kōḷāṟu)}{கோளாறு (kōḷāṟu)}]{\ta{கோளாறு} (kōḷāṟu) — a fault; a malfunction}

\label{lesson:TA-C411-kolaru}



\begin{quote}
Before the new one: say the Tamil for to stroll, then the Tamil for to cross over.
\end{quote}

\subsection*{You'll want to know: \ta{கோளாறு}}
\textbf{\ta{கோளாறு}} — \emph{kōḷāṟu} — \textquotedblleft{}a fault; a malfunction\textquotedblright{}.

\begin{quote}
\ta{இயந்திரத்தில்} \ta{ஏதோ} \ta{கோளாறு} \ta{இருக்கிறது}.
\end{quote}

\begin{quote}
\emph{iyantirattil ētō kōḷāṟu irukkiṟatu.}
\end{quote}

\begin{quote}
There is some fault in the engine.
\end{quote}

Of the body too: \ta{வயிற்றுக்} \ta{கோளாறு}, a stomach upset.

\begin{cousinweb}
\textbf{\ta{கோளாறு}} (\emph{kōḷāṟu}), \textquotedblleft{}a fault, a disorder\textquotedblright{}, is an everyday Tamil word whose origin is not settled.
\end{cousinweb}

\begin{grammarlens}[title={What you've built}]
One more word for getting around.
\end{grammarlens}

\subsection*{Your turn}
\begin{itemize}
\raggedright
  \item \emph{Say it:} \emph{kōḷāṟu}
  \item \emph{Say it:} \emph{kōḷāṟu}, once more
  \item \emph{Say it:} \emph{iyantirattil ētō kōḷāṟu irukkiṟatu.} — There is some fault in the engine.
  \item \emph{Recall:} say the Tamil for to stroll, then the Tamil for to cross over, then say \emph{kōḷāṟu} again
\end{itemize}

\begin{tcolorbox}[breakable,colback=teal!4,colframe=teal!35!black,title={Before you move on}]
What does \ta{கோளாறு} mean? (\textquotedblleft{}A fault; a malfunction\textquotedblright{}.) Say it once more.
\end{tcolorbox}
\section[\texorpdfstring{\ta{செயலிழ} (ceyaliḻa)}{செயலிழ (ceyaliḻa)}]{\ta{செயலிழ} (ceyaliḻa) — to stop working; to fail}

\label{lesson:TA-C411-ceyalila}



\begin{quote}
Before the new one: say the Tamil for to cross over, then the Tamil for a fault.
\end{quote}

\subsection*{You'll want to know: \ta{செயலிழ}}
\textbf{\ta{செயலிழ}} — \emph{ceyaliḻa} — \textquotedblleft{}to stop working; to fail\textquotedblright{}.

\begin{quote}
\ta{நடுவழியில்} \ta{எங்கள்} \ta{காரின்} \ta{இயந்திரம்} \ta{செயலிழந்தது}.
\end{quote}

\begin{quote}
\emph{naṭuvaḻiyil eṅkaḷ kāriṉ iyantiram ceyaliḻantatu.}
\end{quote}

\begin{quote}
Halfway, the engine of our car stopped working.
\end{quote}

Past (it …): \ta{செயலிழந்தது}, \emph{ceyaliḻantatu}.

\ta{செயல்}, an action, + \ta{இழ}, to lose: to lose the power to act. Said of machines and of organs of the body.

\begin{cousinweb}
\textbf{\ta{செயலிழ}} (\emph{ceyaliḻa}) is built from \ta{செயல்}, \textquotedblleft{}an action\textquotedblright{}, the noun of \ta{செய்}, \textquotedblleft{}to do\textquotedblright{} (DEDR 1957), and \ta{இழ}, \textquotedblleft{}to lose\textquotedblright{} (DEDR 501): to lose the power to act. Like \ta{இழ}, it makes its past in -\ta{ந்த}-.
\end{cousinweb}

\begin{grammarlens}[title={What you've built}]
One more word for getting around.
\end{grammarlens}

\subsection*{Your turn}
\begin{itemize}
\raggedright
  \item \emph{Say it:} \emph{ceyaliḻa}
  \item \emph{Say it:} \emph{ceyaliḻa}, once more
  \item \emph{Say it:} \emph{naṭuvaḻiyil eṅkaḷ kāriṉ iyantiram ceyaliḻantatu.} — Halfway, the engine of our car stopped working.
  \item \emph{Recall:} say the Tamil for to cross over, then the Tamil for a fault, then say \emph{ceyaliḻa} again
\end{itemize}

\begin{tcolorbox}[breakable,colback=teal!4,colframe=teal!35!black,title={Before you move on}]
What does \ta{செயலிழ} mean? (\textquotedblleft{}To stop working; to fail\textquotedblright{}.) Say it once more.
\end{tcolorbox}
\section[\texorpdfstring{\ta{நசுங்கு} (nacuṅku)}{நசுங்கு (nacuṅku)}]{\ta{நசுங்கு} (nacuṅku) — to get dented; to be crushed}

\label{lesson:TA-C411-nacunku}



\begin{quote}
Before the new one: say the Tamil for a fault, then the Tamil for to stop working.
\end{quote}

\subsection*{You'll want to know: \ta{நசுங்கு}}
\textbf{\ta{நசுங்கு}} — \emph{nacuṅku} — \textquotedblleft{}to get dented; to be crushed\textquotedblright{}.

\begin{quote}
\ta{விபத்தில்} \ta{காரின்} \ta{கதவு} \ta{நசுங்கியது}.
\end{quote}

\begin{quote}
\emph{vipattil kāriṉ katavu nacuṅkiyatu.}
\end{quote}

\begin{quote}
The car's door got dented in the accident.
\end{quote}

Past (it …): \ta{நசுங்கியது}, \emph{nacuṅkiyatu}.

Of fruit or a can: to be squashed.

\begin{cousinweb}
\textbf{\ta{நசுங்கு}} (\emph{nacuṅku}) is native Dravidian (DEDR 3574), \textquotedblleft{}to be mashed, to be crushed\textquotedblright{}.
\end{cousinweb}

\begin{grammarlens}[title={What you've built}]
One more word for getting around.
\end{grammarlens}

\subsection*{Your turn}
\begin{itemize}
\raggedright
  \item \emph{Say it:} \emph{nacuṅku}
  \item \emph{Say it:} \emph{nacuṅku}, once more
  \item \emph{Say it:} \emph{vipattil kāriṉ katavu nacuṅkiyatu.} — The car's door got dented in the accident.
  \item \emph{Recall:} say the Tamil for a fault, then the Tamil for to stop working, then say \emph{nacuṅku} again
\end{itemize}

\begin{tcolorbox}[breakable,colback=teal!4,colframe=teal!35!black,title={Before you move on}]
What does \ta{நசுங்கு} mean? (\textquotedblleft{}To get dented; to be crushed\textquotedblright{}.) Say it once more.
\end{tcolorbox}
\section[\texorpdfstring{\ta{உரசு} (uracu)}{உரசு (uracu)}]{\ta{உரசு} (uracu) — to scrape; to rub against}

\label{lesson:TA-C411-uracu}



\begin{quote}
Before the new one: say the Tamil for to stop working, then the Tamil for to get dented.
\end{quote}

\subsection*{You'll want to know: \ta{உரசு}}
\textbf{\ta{உரசு}} — \emph{uracu} — \textquotedblleft{}to scrape; to rub against\textquotedblright{}.

\begin{quote}
\ta{குறுகிய} \ta{தெருவில்} \ta{என்} \ta{கார்} \ta{சுவரில்} \ta{உரசியது}.
\end{quote}

\begin{quote}
\emph{kuṟukiya teruvil eṉ kār cuvaril uraciyatu.}
\end{quote}

\begin{quote}
In the narrow street my car scraped the wall.
\end{quote}

Past (I …): \ta{உரசினேன்}, \emph{uraciṉēṉ}.

What is rubbed against takes -\ta{இல்}: \ta{சுவரில்} \ta{உரசு}.

\begin{cousinweb}
\textbf{\ta{உரசு}} (\emph{uracu}) is native Dravidian (DEDR 665), \textquotedblleft{}to rub against\textquotedblright{}.
\end{cousinweb}

\begin{grammarlens}[title={What you've built}]
One more word for getting around.
\end{grammarlens}

\subsection*{Your turn}
\begin{itemize}
\raggedright
  \item \emph{Say it:} \emph{uracu}
  \item \emph{Say it:} \emph{uracu}, once more
  \item \emph{Say it:} \emph{kuṟukiya teruvil eṉ kār cuvaril uraciyatu.} — In the narrow street my car scraped the wall.
  \item \emph{Recall:} say the Tamil for to stop working, then the Tamil for to get dented, then say \emph{uracu} again
\end{itemize}

\begin{tcolorbox}[breakable,colback=teal!4,colframe=teal!35!black,title={Before you move on}]
What does \ta{உரசு} mean? (\textquotedblleft{}To scrape; to rub against\textquotedblright{}.) Say it once more.
\end{tcolorbox}
\section[\texorpdfstring{\ta{முட்டு} (muṭṭu)}{முட்டு (muṭṭu)}]{\ta{முட்டு} (muṭṭu) — to bump; to butt}

\label{lesson:TA-C411-muttu}



\begin{quote}
Before the new one: say the Tamil for to get dented, then the Tamil for to scrape.
\end{quote}

\subsection*{You'll want to know: \ta{முட்டு}}
\textbf{\ta{முட்டு}} — \emph{muṭṭu} — \textquotedblleft{}to bump; to butt\textquotedblright{}.

\begin{quote}
\ta{இருட்டில்} \ta{நான்} \ta{கதவில்} \ta{முட்டினேன்}.
\end{quote}

\begin{quote}
\emph{iruṭṭil nāṉ katavil muṭṭiṉēṉ.}
\end{quote}

\begin{quote}
In the dark I bumped into the door.
\end{quote}

Past (I …): \ta{முட்டினேன்}, \emph{muṭṭiṉēṉ}.

Of a bull or a goat: to butt.

\begin{cousinweb}
\textbf{\ta{முட்டு}} (\emph{muṭṭu}) is native Dravidian (DEDR 4932), \textquotedblleft{}to dash against, to butt\textquotedblright{}.
\end{cousinweb}

\begin{grammarlens}[title={What you've built}]
One more word for getting around.
\end{grammarlens}

\subsection*{Your turn}
\begin{itemize}
\raggedright
  \item \emph{Say it:} \emph{muṭṭu}
  \item \emph{Say it:} \emph{muṭṭu}, once more
  \item \emph{Say it:} \emph{iruṭṭil nāṉ katavil muṭṭiṉēṉ.} — In the dark I bumped into the door.
  \item \emph{Recall:} say the Tamil for to get dented, then the Tamil for to scrape, then say \emph{muṭṭu} again
\end{itemize}

\begin{tcolorbox}[breakable,colback=teal!4,colframe=teal!35!black,title={Before you move on}]
What does \ta{முட்டு} mean? (\textquotedblleft{}To bump; to butt\textquotedblright{}.) Say it once more.
\end{tcolorbox}
